%% T5: noise the method never saw, same frozen configuration. Filled from outputs/paper/noise_kinds.json
\newcommand{\noiserow}[1]{#1 & \pending & \pending & \pending & \pending & \pending & \pending \\}
\begin{table}[t]
  \centering\small
  \caption{\textbf{T5: noise the method never saw.} Late fusion, AP$_{0.7}$, 5 paired seeds, same frozen configuration in every row. U: uncorrected, F: FreeAlign, A: AlignFormer.}
  \begin{tabular}{@{}l*{3}{c}|*{3}{c}@{}}
    \toprule
    & \multicolumn{3}{c|}{OPV2V} & \multicolumn{3}{c}{V2X-Real} \\
    Noise & U & F & A & U & F & A \\
    \midrule
    \noiserow{Benchmark}\noiserow{Slow drift}\noiserow{Constant bias}\noiserow{Mixed quality}\noiserow{Beyond training}
    \bottomrule
  \end{tabular}
  \why{\textbf{Why this table:} the benchmark's noise is independent per frame, which flatters pooling over time, while real GNSS errors persist. This table tests our promise that the method knows nothing about the noise: same settings, noise it never saw. \textbf{How to read it.} Each row stresses one part. Slow drift and constant bias: pooling. Mixed quality: online Abstain, which assumes similar quality across CAVs. Beyond training: the learned Propose. Read A against F in each row, and A against the benchmark row. \textbf{What we conclude.} If A stays above F in every row: the no-knowledge claim holds beyond the benchmark. A row that falls below F names a real limit, and we report it as one.}
\end{table}
